% Problem-section file following ../_template.tex.
\section{Collective cost of tensor-power state preparation}
\label{sec:problem-16}

\paragraph{Problem.}
Determine the asymptotic weighted circuit cost of preparing tensor powers of a
known pure state, and characterize when collective preparation is cheaper per
copy than independent preparation.  On an arbitrary qubit register, allow
Pauli-product rotations $e^{i\phi P}$ with $\phi\in\mathbb{R}$ and
$P\in\{I,X,Y,Z\}^{\otimes q}$, assigning each such gate the cost $|\phi|$.
For a known $m$-qubit state $\lvert\psi\rangle$, let
$U(\boldsymbol\phi,\boldsymbol P):=\prod_{j=1}^r e^{i\phi_jP_j}$ for a finite
sequence of angles and Pauli products.  Define the phase-independent exact
$n$-copy cost by
\begin{equation}
  C_n(\psi)
  :=\inf\left\{
    \sum_{j=1}^{r}|\phi_j|:
    \begin{array}{l}
      r\in\mathbb N_0,\ \phi_j\in\mathbb R,\ 
      P_j\in\{I,X,Y,Z\}^{\otimes mn}\ (1\le j\le r),\\
      U(\boldsymbol\phi,\boldsymbol P)
      \lvert0^{mn}\rangle\!\langle0^{mn}\rvert
      U(\boldsymbol\phi,\boldsymbol P)^\dagger
      =(\lvert\psi\rangle\!\langle\psi\rvert)^{\otimes n}
    \end{array}
  \right\}.
  \label{eq:p16-exact-preparation-cost}
\end{equation}
Determine the growth of Eq.~\eqref{eq:p16-exact-preparation-cost} with $n$
and characterize the states for which the regularized cost
\begin{equation}
  C_\infty(\psi)
  :=\lim_{n\to\infty}\frac{C_n(\psi)}{n}
  =\inf_{n\ge1}\frac{C_n(\psi)}{n}
  \label{eq:p16-regularized-cost}
\end{equation}
satisfies $C_\infty(\psi)<C_1(\psi)$.  For an approximation tolerance
$\varepsilon\ge0$, also determine the scaling of
\begin{equation}
  C_n^\varepsilon(\psi)
  :=\inf_U\left\{
    \operatorname{cost}(U):
    \frac12\left\|
      U\lvert0^{mn}\rangle\!\langle0^{mn}\rvert U^\dagger
      -(\lvert\psi\rangle\!\langle\psi\rvert)^{\otimes n}
    \right\|_1\le\varepsilon
  \right\},
  \label{eq:p16-approximate-preparation-cost}
\end{equation}
where $U$ ranges over all finite circuits of Pauli-product rotations on the
$mn$-qubit register and $\operatorname{cost}(U)$ is the corresponding sum of
absolute rotation angles, as in
Eq.~\eqref{eq:p16-exact-preparation-cost}.  The question for
Eq.~\eqref{eq:p16-approximate-preparation-cost} includes fixed and vanishing
error sequences.

\subsection*{Status}

Unsolved

\subsection*{Source}

The exact and approximate tensor-power preparation questions are posed in the
open-problem collection of Kr\"uger and Werner; Scarani et al. independently
identify collective product-state preparation complexity as an open direction
\sourcecite{ref:p16-krueger-werner}{KW05},
\sourcecite{ref:p16-scarani-et-al}{SIG+05}.

\subsection*{Progress}

\begin{itemize}
  \item Independent preparation gives
  $C_n(\psi)\le nC_1(\psi)$.  Concatenation gives
  $C_{n+k}(\psi)\le C_n(\psi)+C_k(\psi)$, so Fekete's lemma establishes the
  equality in Eq.~\eqref{eq:p16-regularized-cost}.  These elementary bounds do
  not decide whether the inequality can be strict.

  \item Plesch and Brukner gave nearly optimal worst-case state-synthesis
  circuits on a fixed register.  Their discrete local-gate count neither uses
  the weighted Pauli-string metric in
  Eq.~\eqref{eq:p16-exact-preparation-cost} nor exploits a tensor-power promise
  \sourcecite{ref:p16-plesch-brukner}{PB11}.

  \item Mora and Briegel related precision-dependent quantum-state
  algorithmic complexity to entanglement for several state families.  Their
  framework does not determine the direct-product scaling in
  Eq.~\eqref{eq:p16-approximate-preparation-cost}
  \sourcecite{ref:p16-mora-briegel}{MB05}.
\end{itemize}

\subsection*{References}

\begin{enumerate}[label={},leftmargin=4.8em,labelsep=0.5em]
  \footnotesize
  \item[\textup{[PB11]}]\label{ref:p16-plesch-brukner}
  M. Plesch and Č. Brukner,
  ``Quantum-State Preparation with Universal Gate Decompositions,''
  \emph{Physical Review A} \textbf{83}, 032302 (2011).
  \href{https://doi.org/10.1103/PhysRevA.83.032302}%
       {doi:10.1103/PhysRevA.83.032302};
  \href{https://arxiv.org/abs/1003.5760}{arXiv:1003.5760}.

  \item[\textup{[MB05]}]\label{ref:p16-mora-briegel}
  C. Mora and H. J. Briegel,
  ``Algorithmic Complexity and Entanglement of Quantum States,''
  \emph{Physical Review Letters} \textbf{95}, 200503 (2005).
  \href{https://doi.org/10.1103/PhysRevLett.95.200503}%
       {doi:10.1103/PhysRevLett.95.200503};
  \href{https://arxiv.org/abs/quant-ph/0505200}%
       {arXiv:quant-ph/0505200}.

  \item[\textup{[SIG+05]}]\label{ref:p16-scarani-et-al}
  V. Scarani, S. Iblisdir, N. Gisin, and A. Acín,
  ``Quantum Cloning,''
  \emph{Reviews of Modern Physics} \textbf{77}, 1225--1256 (2005).
  \href{https://doi.org/10.1103/RevModPhys.77.1225}%
       {doi:10.1103/RevModPhys.77.1225};
  \href{https://arxiv.org/abs/quant-ph/0511088}%
       {arXiv:quant-ph/0511088}.

  \item[\textup{[KW05]}]\label{ref:p16-krueger-werner}
  O. Krüger and R. F. Werner (eds.),
  ``Some Open Problems in Quantum Information Theory,''
  arXiv:quant-ph/0504166 (2005).
  \href{https://doi.org/10.48550/arXiv.quant-ph/0504166}%
       {doi:10.48550/arXiv.quant-ph/0504166};
  \href{https://arxiv.org/abs/quant-ph/0504166}%
       {arXiv:quant-ph/0504166}.
\end{enumerate}

\subsection*{Comment}

The exact cost model and its approximate variant are posed in the source
collection \sourcecite{ref:p16-krueger-werner}{KW05}; a cloning review also
identifies product-preparation complexity as an open direction
\sourcecite{ref:p16-scarani-et-al}{SIG+05}.  The state is known, so the problem
is not universal cloning.  The unresolved quantity is the collective saving
in Eqs.~\eqref{eq:p16-regularized-cost} and
\eqref{eq:p16-approximate-preparation-cost} for this specific continuous gate
metric.

\subsection*{Tag}

Quantum state preparation; Quantum circuit complexity; Qubit systems

\subsection*{ID}

\texttt{op\_56cf60cdecde44f7}
